\documentclass[11pt,a4paper]{article}
\usepackage[utf8]{inputenc}
\usepackage[spanish]{babel}
\usepackage{amsmath}
\usepackage{amsfonts}
\usepackage{amssymb}
\usepackage{geometry}
\usepackage{enumerate}
\geometry{top=2.5cm, bottom=2.5cm, left=2.5cm, right=2.5cm}

\begin{document}

\begin{center}
    \textbf{\Large FÍSICA BÁSICA} \\[0.2cm]
    \textbf{\large Práctica: Cinemática y Leyes de Newton} \\[0.1cm]
    \small Autor: Eurípides Herasme Medina / derechos reservados \\[0.4cm]
\end{center}

\hr
\section*{I. Preguntas Teóricas (Selecciona la respuesta correcta)}

\begin{enumerate}
    \item \textbf{La fuerza es una magnitud física de carácter:}
    \begin{enumerate}[a)]
        \item Escalar.
        \item Vectorial.
        \item Adimensional.
        \item Tensorial de segundo orden.
    \end{enumerate}

    \item \textbf{De acuerdo con la Primera Ley de Newton, si la fuerza neta que actúa sobre un cuerpo es nula, el cuerpo necesariamente:}
    \begin{enumerate}[a)]
        \item Debe estar en reposo absoluto.
        \item Permanecerá en reposo o en movimiento rectilíneo uniforme.
        \item Experimentará un aumento constante de su velocidad.
        \item Cambiará continuamente su dirección de movimiento.
    \end{enumerate}

    \item \textbf{La Segunda Ley de Newton establece matemáticamente que la aceleración de un cuerpo es:}
    \begin{enumerate}[a)]
        \item Directamente proporcional a su masa e inversamente proporcional a la fuerza neta.
        \item Directamente proporcional a la fuerza neta e inversamente proporcional a su masa.
        \item Independiente de la masa y proporcional al peso.
        \item Proporcional al cuadrado de la fuerza neta aplicada.
    \end{enumerate}

    \item \textbf{De acuerdo con la Tercera Ley de Newton (Acción-Reacción), si un cuerpo A ejerce una fuerza $\vec{F}_{AB}$ sobre un cuerpo B, entonces:}
    \begin{enumerate}[a)]
        \item El cuerpo B ejerce sobre A una fuerza de igual magnitud, misma dirección y sentido contrario ($\vec{F}_{BA} = -\vec{F}_{AB}$).
        \item El cuerpo B ejerce sobre A una fuerza de mayor magnitud en el mismo sentido.
        \item La fuerza del cuerpo B sobre A es nula porque A inició el contacto.
        \item El cuerpo B ejerce sobre A una fuerza perpendicular a la original.
    \end{enumerate}

    \item \textbf{La propiedad física de un cuerpo que mide cuantitativamente su inercia (resistencia al cambio en su estado de movimiento) es la:}
    \begin{enumerate}[a)]
        \item Aceleración.
        \item Fuerza normal.
        \item Masa.
        \item Fricción.
    \end{enumerate}

    \item \textbf{En un gráfico de posición en función del tiempo ($x$ vs $t$), la velocidad instantánea en un tiempo determinado equivale a:}
    \begin{enumerate}[a)]
        \item El área total bajo la curva en ese intervalo.
        \item La pendiente de la recta secante que corta la curva.
        \item La pendiente de la recta tangente a la curva en ese punto.
        \item La longitud del vector de aceleración.
    \end{enumerate}

    \item \textbf{Con respecto a la fuerza de fricción entre dos superficies secas en contacto, se cumple que:}
    \begin{enumerate}[a)]
        \item El coeficiente de fricción cinética ($\mu_c$) es siempre menor que el coeficiente de fricción estática ($\mu_e$).
        \item La fuerza de fricción estática tiene un valor constante e invariable.
        \item La fuerza de fricción es siempre perpendicular a la superficie de contacto.
        \item El coeficiente de fricción depende del área de contacto visible.
    \end{enumerate}

    \item \textbf{En un Movimiento Circular Uniforme (MCU), aunque la rapidez es constante, la velocidad cambia de dirección, lo que implica una aceleración centrípeta dirigida hacia:}
    \begin{enumerate}[a)]
        \item La dirección tangente a la trayectoria circular.
        \item El centro de la trayectoria circular.
        \item El sentido opuesto del movimiento lineal.
        \item El exterior de la circunferencia en dirección radial.
    \end{enumerate}

    \item \textbf{Despreciando la resistencia del aire, el ángulo de disparo respecto a la horizontal que produce el máximo alcance horizontal para un proyectil es:}
    \begin{enumerate}[a)]
        \item $30^\circ$
        \item $45^\circ$
        \item $60^\circ$
        \item $90^\circ$
    \end{enumerate}

    \item \textbf{A diferencia de la masa, el peso de un objeto:}
    \begin{enumerate}[a)]
        \item Es una propiedad intrínseca que no varía con la gravedad local.
        \item Es la fuerza de atracción gravitatoria que ejerce la Tierra sobre él ($w = m \cdot g$).
        \item Se mide en kilogramos en el Sistema Internacional.
        \item Es una cantidad escalar pura sin dirección ni sentido.
    \end{enumerate}
\end{enumerate}

\hr
\section*{II. Problemas Prácticos (Selecciona la respuesta correcta)}

\begin{enumerate}
    \item \textbf{Una partícula recibe la acción de dos fuerzas coplanares: $\vec{F}_1 = (1.50 \text{ N}, 1.60 \text{ N})$ y $\vec{F}_2 = (2.50 \text{ N}, 1.40 \text{ N})$. Determine las componentes de la fuerza neta:}
    \begin{enumerate}[a)]
        \item $(3.00 \text{ N}, 4.00 \text{ N})$
        \item $(4.00 \text{ N}, 3.00 \text{ N})$
        \item $(1.00 \text{ N}, 0.20 \text{ N})$
        \item $(4.00 \text{ N}, 4.00 \text{ N})$
    \end{enumerate}

    \item \textbf{Determine la magnitud de la fuerza neta sobre una partícula que recibe la acción de las fuerzas $\vec{F}_1 = (5.00 \text{ N}, 4.00 \text{ N})$ y $\vec{F}_2 = (3.00 \text{ N}, 2.00 \text{ N})$:}
    \begin{enumerate}[a)]
        \item $14.0 \text{ N}$
        \item $10.0 \text{ N}$
        \item $8.00 \text{ N}$
        \item $6.00 \text{ N}$
    \end{enumerate}

    \item \textbf{Un bloque en una superficie horizontal recibe tres fuerzas paralelas al eje $x$: $F_{1x} = -5.00 \text{ N}$ (izquierda), $F_{2x} = +9.00 \text{ N}$ (derecha) y $F_{3x} = -7.00 \text{ N}$ (izquierda). Determine el valor de la fuerza neta horizontal:}
    \begin{enumerate}[a)]
        \item $+3.00 \text{ N}$
        \item $-3.00 \text{ N}$
        \item $+11.0 \text{ N}$
        \item $-21.0 \text{ N}$
    \end{enumerate}

    \item \textbf{Calcule la magnitud de la fuerza neta que actúa sobre un bloque de masa $m = 4.00 \text{ kg}$ si este se mueve con una aceleración constante de $1.20 \text{ m/s}^2$:}
    \begin{enumerate}[a)]
        \item $3.33 \text{ N}$
        \item $4.80 \text{ N}$
        \item $5.20 \text{ N}$
        \item $0.30 \text{ N}$
    \end{enumerate}

    \item \textbf{Determine la masa de un cuerpo que se mueve con una aceleración de $2.00 \text{ m/s}^2$ sabiendo que la fuerza neta que actúa sobre él es $\vec{F}_{\text{neta}} = (5.00 \text{ N}, 7.48 \text{ N})$:}
    \begin{enumerate}[a)]
        \item $2.50 \text{ kg}$
        \item $4.50 \text{ kg}$
        \item $6.24 \text{ kg}$
        \item $9.00 \text{ kg}$
    \end{enumerate}

    \item \textbf{Un bloque de masa $m = 4.00 \text{ kg}$ es sometido a dos fuerzas horizontales: $F_{1x} = -8.00 \text{ N}$ (hacia la izquierda) y $F_{2x} = +10.0 \text{ N}$ (hacia la derecha). Determine la componente de su aceleración $a_x$:}
    \begin{enumerate}[a)]
        \item $+0.500 \text{ m/s}^2$
        \item $-0.500 \text{ m/s}^2$
        \item $+4.50 \text{ m/s}^2$
        \item $+2.00 \text{ m/s}^2$
    \end{enumerate}

    \item \textbf{Un cuerpo de masa $m = 3.00 \text{ kg}$ acelera horizontalmente a razón de $a_x = +0.800 \text{ m/s}^2$ (hacia la derecha). Si actúan dos fuerzas sobre el eje $x$ y una de ellas es $F_{2x} = +9.00 \text{ N}$, determine la componente de la otra fuerza $F_{1x}$:}
    \begin{enumerate}[a)]
        \item $+6.60 \text{ N}$
        \item $-6.60 \text{ N}$
        \item $-11.4 \text{ N}$
        \item $-2.40 \text{ N}$
    \end{enumerate}

    \item \textbf{Un automóvil que viaja en línea recta a $50.0 \text{ km/h}$ ($13.9 \text{ m/s}$) pisa los frenos durante $3.00 \text{ s}$ reduciendo su velocidad de forma uniforme hasta $18.0 \text{ km/h}$ ($5.00 \text{ m/s}$). Calcule la distancia recorrida ($\Delta x$) durante el frenado:}
    \begin{enumerate}[a)]
        \item $15.0 \text{ m}$
        \item $28.4 \text{ m}$
        \item $40.0 \text{ m}$
        \item $56.7 \text{ m}$
    \end{enumerate}

    \item \textbf{Una caja se desliza por una rampa inclinada con una aceleración constante de $2.00 \text{ m/s}^2$ partiendo con una velocidad inicial de $1.50 \text{ m/s}$. Calcule el desplazamiento de la caja al cabo de $3.00 \text{ s}$:}
    \begin{enumerate}[a)]
        \item $9.00 \text{ m}$
        \item $13.5 \text{ m}$
        \item $18.0 \text{ m}$
        \item $22.5 \text{ m}$
    \end{enumerate}

    \item \textbf{Una bola de golf se desliza por el césped con una velocidad inicial de $6.00 \text{ m/s}$ y llega al hoyo a $2.00 \text{ m/s}$ después de recorrer una distancia de $10.0 \text{ m}$ en línea recta. Determine su aceleración constante:}
    \begin{enumerate}[a)]
        \item $-1.60 \text{ m/s}^2$
        \item $+1.60 \text{ m/s}^2$
        \item $-3.20 \text{ m/s}^2$
        \item $-0.80 \text{ m/s}^2$
    \end{enumerate}

    \item \textbf{Un estudiante deja caer un tomate (desde el reposo) desde la azotea de un edificio y mide que tarda exactamente $1.60 \text{ s}$ en golpear el suelo. Despreciando la resistencia del aire y usando $g = 9.80 \text{ m/s}^2$, calcule la altura del edificio:}
    \begin{enumerate}[a)]
        \item $15.7 \text{ m}$
        \item $12.5 \text{ m}$
        \item $25.1 \text{ m}$
        \item $7.84 \text{ m}$
    \end{enumerate}

    \item \textbf{Determine la magnitud de la fuerza de fricción cinética que actúa sobre un bloque que se desliza por el suelo horizontal, si la fuerza normal de la superficie sobre él es $n = 8.00 \text{ N}$ y el coeficiente de fricción cinético es $\mu_c = 0.200$:}
    \begin{enumerate}[a)]
        \item $40.0 \text{ N}$
        \item $1.60 \text{ N}$
        \item $0.025 \text{ N}$
        \item $9.80 \text{ N}$
    \end{enumerate}

    \item \textbf{Un cuerpo pesa $14.7 \text{ N}$ en la Tierra ($g = 9.80 \text{ m/s}^2$). ¿Cuánto pesará este mismo cuerpo en la superficie de la Luna, donde la aceleración de la gravedad es de $1.62 \text{ m/s}^2$?:}
    \begin{enumerate}[a)]
        \item $14.7 \text{ N}$
        \item $2.43 \text{ N}$
        \item $23.8 \text{ N}$
        \item $1.50 \text{ N}$
    \end{enumerate}

    \item \textbf{Una cápsula de pasajeros en una rueda de la fortuna de $150 \text{ m}$ de diámetro gira a rapidez constante y le toma exactamente $30.0 \text{ minutos}$ ($1800 \text{ s}$) completar una vuelta entera. Calcule su rapidez lineal:}
    \begin{enumerate}[a)]
        \item $0.524 \text{ m/s}$
        \item $0.262 \text{ m/s}$
        \item $1.50 \text{ m/s}$
        \item $5.00 \text{ m/s}$
    \end{enumerate}

    \item \textbf{Un estudiante mide que las aspas de un abanico de techo realizan exactamente $40 \text{ vueltas}$ en un intervalo de $2.50 \text{ s}$. Calcule la velocidad angular $\omega$ del abanico en rad/s (utilice $\pi \approx 3.14$):**}
    \begin{enumerate}[a)]
        \item $16.0 \text{ rad/s}$
        \item $100 \text{ rad/s}$
        \item $50.2 \text{ rad/s}$
        \item $251 \text{ rad/s}$
    \end{enumerate}
\end{enumerate}

\hr
\section*{Clave de Respuestas (Para el docente)}

\subsection*{Sección I: Preguntas Teóricas}
\begin{enumerate}
    \item \textbf{b)} Vectorial. (La fuerza tiene dirección, sentido y magnitud).
    \item \textbf{b)} Permanecerá en reposo o en movimiento rectilíneo uniforme. (Primera ley de Newton / Inercia).
    \item \textbf{b)} Directamente proporcional a la fuerza neta e inversamente proporcional a su masa. ($\vec{a} = \vec{F}_{\text{neta}}/m$).
    \item \textbf{a)} El cuerpo B ejerce sobre A una fuerza de igual magnitud, misma dirección y sentido contrario ($\vec{F}_{BA} = -\vec{F}_{AB}$). (Tercera ley de Newton).
    \item \textbf{c)} Masa. (La masa es la medida cuantitativa de la inercia).
    \item \textbf{c)} La pendiente de la recta tangente a la curva en ese punto. ($v_x(t) = dx/dt$).
    \item \textbf{a)} El coeficiente de fricción cinética ($\mu_c$) es siempre menor que el coeficiente de fricción estática ($\mu_e$).
    \item \textbf{b)} El centro de la trayectoria circular. (La aceleración centrípeta es radial hacia el centro).
    \item \textbf{b)} $45^\circ$. (Produce el máximo término en la ecuación de alcance horizontal).
    \item \textbf{b)} Es la fuerza de atracción gravitatoria que ejerce la Tierra sobre él ($w = m \cdot g$). (El peso varía con la gravedad local).
\end{enumerate}

\subsection*{Sección II: Problemas Prácticos}
\begin{enumerate}
    \item \textbf{b)} $(4.00 \text{ N}, 3.00 \text{ N})$. Sumando las componentes: $F_x = 1.50 + 2.50 = 4.00 \text{ N}$; $F_y = 1.60 + 1.40 = 3.00 \text{ N}$.
    \item \textbf{b)} $10.0 \text{ N}$. $\vec{F}_{\text{neta}} = (8.00\text{ N}, 6.00\text{ N})$, y luego su magnitud: $F = \sqrt{8^2 + 6^2} = 10.0 \text{ N}$.
    \item \textbf{b)} $-3.00 \text{ N}$. $F_{\text{neta-}x} = -5.00 + 9.00 - 7.00 = -3.00 \text{ N}$.
    \item \textbf{b)} $4.80 \text{ N}$. Segunda ley: $F = m \cdot a = 4.00 \cdot 1.20 = 4.80 \text{ N}$.
    \item \textbf{b)} $4.50 \text{ kg}$. Fuerza neta: $F = \sqrt{5.00^2 + 7.48^2} = 9.00 \text{ N}$. Masa: $m = F / a = 9.00 / 2.00 = 4.50 \text{ kg}$.
    \item \textbf{a)} $+0.500 \text{ m/s}^2$. $F_{\text{neta-}x} = -8.00 + 10.0 = +2.00 \text{ N}$. $a_x = F_x / m = 2.00 / 4.00 = +0.500 \text{ m/s}^2$.
    \item \textbf{b)} $-6.60 \text{ N}$. $F_x = m \cdot a_x = 3.00 \cdot 0.800 = +2.40 \text{ N}$. $F_{1x} = 2.40 - 9.00 = -6.60 \text{ N}$.
    \item \textbf{b)} $28.4 \text{ m}$. $\Delta x = \frac{v_x + v_{0x}}{2} \cdot t = \frac{13.9 + 5.00}{2} \cdot 3.00 = 28.35 \text{ m} \approx 28.4 \text{ m}$.
    \item \textbf{b)} $13.5 \text{ m}$. $\Delta x = v_{0x} t + \frac{1}{2} a_x t^2 = 1.50(3.00) + 0.5(2.00)(3.00)^2 = 13.5 \text{ m}$.
    \item \textbf{a)} $-1.60 \text{ m/s}^2$. $a_x = \frac{v_x^2 - v_{0x}^2}{2\Delta x} = \frac{2.00^2 - 6.00^2}{2(10.0)} = -1.60 \text{ m/s}^2$.
    \item \textbf{b)} $12.5 \text{ m}$. Altura: $h = |\Delta y| = \frac{1}{2} g t^2 = 0.5(9.80)(1.60)^2 = 12.54 \text{ m} \approx 12.5 \text{ m}$.
    \item \textbf{b)} $1.60 \text{ N}$. Fricción cinética: $f_c = \mu_c \cdot n = 0.200 \cdot 8.00 = 1.60 \text{ N}$.
    \item \textbf{b)} $2.43 \text{ N}$. Masa: $m = w / g = 14.7 / 9.80 = 1.50 \text{ kg}$. Peso Luna: $w_L = m \cdot g_L = 1.50 \cdot 1.62 = 2.43 \text{ N}$.
    \item \textbf{b)} $0.262 \text{ m/s}$. Radio $R = 75 \text{ m}$. $v = \frac{2\pi R}{T} = \frac{2(3.1416)(75)}{1800} \approx 0.262 \text{ m/s}$.
    \item \textbf{b)} $100 \text{ rad/s}$. $f = 40 / 2.50 = 16.0 \text{ Hz}$. $\omega = 2\pi f = 2(3.14)(16.0) = 100.48 \approx 100 \text{ rad/s}$.
\end{enumerate}

\end{document}
